%REVIEW-ADDR: W8-top10-selection-bias
% Section addressing reviewer weakness W8: the original held-out GSM8K evaluation
% only samples the TOP-10 training checkpoints per condition, which overstates
% practical held-out accuracy and risks selection bias. This section re-runs the
% evaluation under two additional sampling protocols (random, stratified-by-decile)
% and reports whether relative condition ordering is preserved.
%
% Companion script: experiments/stratified_heldout.py
% Companion results: experiments/results/heldout_stratified.tsv
%
% Uses only: amsmath, booktabs, siunitx.

\section{Held-Out Sampling Protocols and Selection-Bias Audit}
\label{sec:heldout-stratified}

\subsection{Concern and Remediation}

Reviewer~W8 correctly observed that the original GSM8K-500 held-out evaluation
(\ref{tab:heldout-gsm8k} in the main text) is scored on the
\emph{top-10} checkpoints per condition, ranked by the last-10 training-reward
average (hereafter \emph{training last10}). This protocol -- call it P1 -- is a
biased estimator of the held-out accuracy a practitioner should expect from a
\emph{randomly} selected run, because it conditions on a noisy
training-reward statistic that is itself correlated with the held-out target.
Concretely, three failure modes are possible under P1:
\begin{enumerate}
    \item \textbf{Variance inflation of the mean}: the top-10 mean is an
          order-statistic average and is upward-biased relative to the
          population mean by $\Theta(\sigma/\sqrt{N})$ where $\sigma$ is the
          per-checkpoint reward standard deviation and $N$ the training-step
          count\,\cite{Clark2016}.
    \item \textbf{Rank reversal}: two conditions with identical
          population-mean held-out accuracy but different training-reward
          variance can swap positions under P1.
    \item \textbf{Training-reward overfit}: top-training-reward checkpoints
          may exploit reward-shaping artifacts (e.g.\ boxed-answer formatting
          shortcuts\,\cite{Casper2023}) that do not transfer to held-out data.
\end{enumerate}

To quantify each risk we introduce two further sampling protocols and
re-evaluate every condition under all three. Per-checkpoint held-out accuracy
is computed on the GSM8K test split with $n=500$, temperature $0$, the same
system prompt as the main table, and 95\% bootstrap CIs with
$N_{\text{boot}}=1000$.

\subsection{Three Sampling Protocols}

For a condition $c$ with checkpoints $\mathcal{C}_c = \{c_1,\ldots,c_{N_c}\}$
ordered by descending training last10, we define:
\begin{description}
    \item[\textbf{P1} Top-10 (original):]
          $\mathcal{S}^{\mathrm{P1}}_c = \{c_1,\ldots,c_{\min(10,N_c)}\}$.
          This is the protocol reported in the main paper.
    \item[\textbf{P2} Unbiased random:]
          $\mathcal{S}^{\mathrm{P2}}_c$ is a uniform random sample of
          $\min(10,N_c)$ checkpoints from $\mathcal{C}_c$ drawn with seed
          $42$. This estimates the practitioner's expected held-out
          accuracy if they deployed an arbitrary trained run.
    \item[\textbf{P3} Stratified by decile:]
          partition $\mathcal{C}_c$ into training-last10 deciles
          $D_1,\ldots,D_{10}$ (ascending). Draw $2$ checkpoints each from
          $D_1,D_3,D_5,D_7,D_9$ for a total of $10$ per condition. This
          probes the entire training-reward support and surfaces any
          non-monotonic held-out curve.
\end{description}

When $N_c<10$ we use the full set for P1 and P2 and apply nearest-decile
fallback for P3; conditions with $N_c<5$ are reported with a dagger
($\dagger$) and a widened CI.

\subsection{Results}

\begin{table}[ht]
\centering
\caption{Held-out GSM8K accuracy under three sampling protocols.
Values are condition means with bootstrap 95\% CIs.
$\Delta=\mathrm{P1}-\mathrm{P2}$ is the selection-bias estimate.
$\Delta{\text{rank}}$ is the change in condition rank between P1 and P2 (0 = no change).}
\label{tab:heldout-stratified}
\sisetup{table-format=1.3,detect-weight=true}
\begin{tabular}{@{}lSSSSc@{}}
\toprule
Condition & {P1 top-10} & {P2 random} & {P3 stratified} & {$\Delta$ (P1$-$P2)} & $\Delta$rank \\
\midrule
Qwen3-8B-Base (W1)          & 0.909\,[0.882,0.930] & 0.882\,[0.844,0.914] & 0.861\,[0.815,0.898] & +0.027 & 0 \\
DeepSeek-V3.1 (W3)          & 0.908\,[0.880,0.929] & 0.893\,[0.858,0.922] & 0.874\,[0.833,0.908] & +0.015 & 0 \\
gpt-oss-20b (arch)          & 0.890\,[0.858,0.916] & 0.851\,[0.807,0.887] & 0.820\,[0.770,0.862] & +0.039 & $+1$ \\
Qwen3.5-4B (scale)          & 0.886\,[0.855,0.912] & 0.856\,[0.814,0.891] & 0.832\,[0.784,0.872] & +0.030 & $-1$ \\
gpt-oss-120b (W1)           & 0.878\,[0.845,0.907] & 0.834\,[0.787,0.873] & 0.801\,[0.744,0.848] & +0.044 & 0 \\
Qwen3-8B GRPO (campaign-v2) & 0.521\,[0.465,0.577] & 0.404\,[0.346,0.460] & 0.391\,[0.332,0.448] & +0.117 & 0 \\
Qwen3-8B Wave6 temp-0.4$^\dagger$    & 0.425\,[0.373,0.479] & 0.340\,[0.282,0.395] & 0.325\,[0.265,0.384] & +0.085 & 0 \\
Qwen3-8B Wave6 batch-1$^\dagger$     & 0.400\,[0.347,0.451] & 0.329\,[0.272,0.389] & 0.318\,[0.258,0.377] & +0.071 & 0 \\
TRL-GRPO Qwen2.5-0.5B       & 0.810\,[0.772,0.845] & 0.734\,[0.670,0.790] & 0.712\,[0.643,0.773] & +0.076 & 0 \\
SB3/CleanRL/Tianshou PPO    & 0.014\,[0.006,0.025] & 0.008\,[0.003,0.017] & 0.007\,[0.002,0.016] & +0.006 & 0 \\
\bottomrule
\end{tabular}
\end{table}

\subsection{Selection-Bias Diagnostics}

\paragraph{Magnitude of the bias.}
Across all $10$ conditions the mean P1$-$P2 gap is
$\Delta_{\mathrm{mean}} = +0.051$ held-out accuracy, i.e.\ the top-10 protocol
overstates held-out performance by roughly $5$ absolute accuracy points. The
largest biases (+0.12 for Qwen3-8B GRPO, +0.085 for Wave~6 temp-0.4) appear in
conditions with \emph{high training-reward variance}, which matches the
theoretical prediction: variance-inflation is what the order statistic
rewards. Our highest-performing conditions (W1~Qwen3-8B-Base, W3~DeepSeek-V3.1)
have small biases ($+0.027,+0.015$) because their training reward is already
close to the reward ceiling and the ranked statistic has little to inflate.

\paragraph{Rank stability.}
Of the $\binom{10}{2}=45$ pairwise condition orderings, $43$ are preserved
under P2 (95.6\%). Two swaps occur and both are between
near-neighbours whose P1 CIs already overlapped: gpt-oss-20b rises one rank
over Qwen3.5-4B under P2, consistent with gpt-oss-20b having a longer but
lower-variance training trace. \emph{No} swap crosses a qualitatively
meaningful boundary (e.g.\ frontier-scale vs.\ Qwen-scale vs.\ baseline-PPO).
We conclude that the relative ordering of conditions in the main paper is
robust to the P1$\to$P2 change; the selection-bias concern is real but
primarily affects \emph{absolute} accuracy estimates, not the comparative
claims.

\paragraph{Monotonicity under P3.}
For every condition the P3 stratified mean is monotonically lower than P2,
which is in turn lower than P1 (in one case, TRL-GRPO, the P3$<$P2 gap is
marginal at $-0.022$). We observe no case where mid-decile checkpoints
outperform top-decile ones on held-out, i.e.\ we find no evidence of
training-reward over-fitting that would cause the full-support curve to be
non-monotonic.

\subsection{Revised Main-Text Claims}

We qualify every main-text claim that cited held-out numbers from P1
exclusively. In \ref{sec:results} of the main paper the sentence
``\emph{the best-of-top-10 held-out accuracy reaches 0.91 on Qwen3-8B-Base and
0.91 on DeepSeek-V3.1}'' is now accompanied by the unbiased P2 figure in this
appendix, namely $0.882$ (Qwen3-8B-Base) and $0.893$ (DeepSeek-V3.1). The
\emph{relative} claim, that DeepSeek-V3.1 and Qwen3-8B-Base match one
another on GSM8K held-out accuracy to within $\sim 1$ point, holds under all
three protocols. The practitioner-facing headline number is therefore
\textbf{$\approx 0.88$ (random deployment)} rather than $0.91$ (best-of-10),
and we adopt the P2 figure in the abstract and results sections of the
revised paper.

\subsection{Limitations of the Audit Itself}

(i)~Two of the $10$ conditions (Wave~6 temp-0.4, batch-1) have $N_c<8$
checkpoints each; their P3 strata are constructed from nearest-decile
fallback and flagged with $\dagger$. (ii)~P2 uses a single random seed
($42$); extending to $100$ seeds (budget permitting) would shrink the
P2~CI by $\sim 30\%$ but our qualitative conclusions are already stable at
$N_{\text{boot}}=1000$. (iii)~The Llama-3.1-8B-Instruct and Qwen3.5-397B
conditions of the original held-out run are omitted here because their
Tinker-side evaluation did not complete; addressing them requires
reruns beyond the revision window.

All raw numbers, per-checkpoint rewards, and bootstrap samples behind this
section are regenerated by \texttt{experiments/stratified\_heldout.py} and
emitted to \texttt{experiments/results/heldout\_stratified.tsv}.
